\slajdid{10-s084}
\begin{frame}[label=10-s084]
\gusttekst\setlength{\abovedisplayskip}{3pt}\setlength{\belowdisplayskip}{3pt}\setlength{\abovedisplayshortskip}{2pt}\setlength{\belowdisplayshortskip}{2pt}\begin{columns}[T,onlytextwidth]\column{.29\textwidth}\centering\input{slike/tikz/s082_ecl_diferencijalni_par.tex}\column{.69\textwidth}
Prilikom povećanja ulaznog napona, smanjuje se inverzna polarizacija BE spoja tranzistora T1, i pri nekom ulaznom naponu BE spoj postaje pozitivno polarizovan i mogu da se stvore uslovi da i tranzistor T1 počne da vodi
\[V_I=V_{BE1}-V_{BE2}+V_R\]
\[V_I=V_{\gamma T}-V_{BE}+V_R\]
Pri tom ulaznom naponu još uvek sva struja $I_E$ prolazi kroz tranzistor T2, ali prilikom daljeg povećavanja ulaznog napona počinje da se pojavljuje i struja kroz tranzistor T1. U situaciji kada su oba tranzistora identično polarizovana ulazni napon je
\[V_I=V_{BE}-V_{BE}+V_R=V_R\]
\end{columns}\par\smallskip
Tada kroz svaki tranzistor prolazi polovina struje $I_E$ pa je
\[V_{O2}=V_{O1}=V_{CC}-R_{C2}\frac{I_{E2}}2=V_{CC}-R_{C1}\frac{I_{E2}}2\]
odnosno bira se $R_{C2}=R_{C1}=R_C$ tako da taj uslov bude ispunjen. Daljim porastom ulaznog napona sve više struje prolazi kroz transistor T1, sve manje kroz transistor T2, dok tranzistor T2 ne dođe na ivicu provođenja pri naponu
\[V_I=V_{BE}-V_{\gamma T}+V_R\]

\end{frame}
